% !TEX root = ../main.tex
\chapter{Jobs and long runs}\label{ch:jobs}
Most chains finish in seconds. Some take minutes, and a few would take hours or days on a laptop.
The app handles long work in three ways. Work that fits runs here, on the machine serving the site,
as a \term{job}{job} you can watch. Work too long for this machine is written as a job for the
university's computing cluster, which you send and bring back yourself. And work that waits for a
person, such as a queue of detections to judge, waits in Review. The Jobs page shows all three in
one place.

\begin{plainterms}[Here, the cluster, and review]
Think of a kitchen. Some dishes cook on the hob while you watch. A big roast goes to the baker's
oven down the road: you write out the instructions, carry the tray there, and collect it later.
And some plates wait at the pass for someone to taste them. Jobs is the board on the kitchen wall
that lists all three.
\end{plainterms}

\pagesection{jobs.all}{All jobs}{In progress here, waiting on the cluster, waiting in review}{live}
One page, with three boards stacked in the order you wait on them: \emph{In progress here},
\emph{Waiting on the cluster}, and \emph{Waiting in review}, plus the \ui{Manifest inbox} in a drawer.
Every row comes from the bridge; nothing on the page is example data. The boards re-read
themselves every few seconds while the tab is visible, faster while something is running.

\pagefield{When you'd come here}
\begin{itemize}
\item A training run, a window set, a tree or an import is running and you want to see how far it
  has got, or why it failed.
\item You wrote a job for the cluster and want to know what to copy, or its results are back and
  need importing.
\item You want to know how much is left to judge in Review.
\end{itemize}

\screenshot{pages/jobs/jobs.all--default}{The Jobs page. The chips under the title count what is in
progress here, waiting on the cluster, back to import, and left to judge; each chip scrolls to its
board.}

\pagefield{Doing things here}

\subsection{Watch work running on this machine}\label{sec:jobs-local}
The first board, \emph{In progress here · this bridge's jobs}, lists the bridge's own jobs: chain
runs, model training, window sets and trees, blind-queue work, sweeps and imports. Running jobs come
first, then the newest.
\begin{enumerate}
\item Read each row: its id, what it is doing (the stage), its state with a progress bar when the
  job reports a fraction, when it started, the estimate (for a chain run, from each block's own
  estimator), and how long it took.
\item Narrow the list with the two switches at the right of the board: by state (\ui{all},
  \ui{in progress}, \ui{failed}, \ui{finished}) and by kind.
\item Click a row to open it. A failed job shows its error and the full traceback.
\item Press \ui{Cancel} to stop a running job. A chain run has no \ui{Cancel} here; stop it from its
  own page in Analyse. A job stops at the next point where it checks, not mid-calculation.
\item Press the link at the end of a row to open its result in the workspace that reads it.
\item Press \ui{show older jobs} to list more than the first six.
\end{enumerate}

\screenshot{pages/zzzzzzzzz_aj_jobs/jobs.all--fixup-aj-running-here-is-the-bridges-own-job-table}{In
progress here: the bridge's own job table, with state, stage and timing for each job.}

\subsection{Send a job to the cluster}\label{sec:jobs-cluster}
The second board, \emph{Waiting on the cluster · scripts the site wrote}, lists every job the site
has written for the cluster. The site does not write them here; you write one from the page whose
work it is (\cref{sec:jobs-sending}). Then:
\begin{enumerate}
\item Click the job's row. It lists \emph{what to copy to the cluster}, with each file's size and
  its path relative to the repository, and the command to run there, from the repository root, such
  as \code{sbatch HPC/Training/generated/<job>/<job>.sh}.
\item Copy those files to the same paths in your copy of the repository on the cluster, and run
  the command. The site never logs in to the cluster or submits anything itself.
\item Wait. The row says \emph{written · waiting for results}. A seed search that has sent back a
  checkpoint says \emph{running on the cluster · checkpoint back}, with how many channels are done.
\item When the job has finished, copy its results back to the place the row names: a job folder's
  \code{out/} folder, or a seed search's result file. Press \ui{Look again}; the row now says
  \emph{results back · not imported}.
\item Press \ui{Import results} (or \ui{Import result} for a seed search). The results are checked
  before anything is stored. The row then says \emph{results imported}, with the run, and its link
  opens it; or \emph{import refused}, with the reason.
\end{enumerate}
A row whose results cannot come back through the site (Discovery's chain export, Models' paired
training) says so instead of offering an import.

\screenshot{pages/jobs/jobs.all--cluster-row-opened}{An opened cluster row: what to copy, how big it
is, and the sbatch command to run on the cluster.}

\subsection{Import a returned job folder}\label{sec:jobs-inbox}
\begin{enumerate}
\item Press \ui{Manifest inbox · import results} at the top of the page. A drawer opens.
\item Type the path of the job folder the results were copied back into, or press \ui{use this
  folder} beside one the site already sees as back.
\item Press \ui{Import}. The drawer follows the import and says \emph{Results imported},
  \emph{Already imported}, or \emph{Import refused — nothing was recorded} with the reason.
\item When it is imported, press the button that opens the result in Models \textrightarrow{}
  Results or in Analyse.
\end{enumerate}
The drawer also lists recent import attempts. A seed search's result is a file, not a folder: import
it from its row on the board.

\screenshot{pages/zzzzzzzzz_aj_jobs/jobs.all--fixup-aj-the-manifest-inbox-refuses-a-folder-with-the-reason}{The
Manifest inbox refuses a folder whose results do not match the job it was written as, and says
why. Nothing is stored.}

\subsection{See what waits in review}\label{sec:jobs-review}
The third board, \emph{Waiting in review · open queues}, lists every open
\term{review queue}{review-queue}: what it judges, how many are judged out of how many, how many are
left, and the pace so far with a time left. Press \ui{Judge} to open the queue in Review
(\apppage{review.inspector}), or \ui{Open Review} for the workspace.

\begin{watchout}[A job written but not sent is not running]
``Waiting on the cluster'' means the site wrote the job. It does not know whether you sent it. A
row stays \emph{written} until results appear where it expects them.
\end{watchout}

\pagefield{Where it leads}
A finished local job links to its result: Analyse for a chain or a tree, Models for training
(\apppage{models.results}). An imported cluster job opens in Models \textrightarrow{} Results,
Analyse or Discovery (\apppage{discovery.seed}), depending on what wrote it. A review queue opens in
Review (\apppage{review.inspector}).

\begin{underthehood}[Jobs page]
Client: \code{webui/client/src/jobs/AllPage.tsx} and \code{LiveJobs.tsx}, reading through
\code{webui/client/src/api/jobs.ts}. Local jobs: \code{GET /api/jobs}
(\code{webui/server/jobs.py}, \code{JobManager.list_jobs}), cancel \code{POST
/api/jobs/{id}/cancel}. Cluster jobs: \code{GET /api/hpc/cluster}
(\code{webui/server/jobs_routes.py}), which joins the job folders
(\code{Working.hpc.job_export.list_exported_jobs}), Discovery's seed searches sent to the cluster,
and the flat recipe-and-script pairs. Inbox: \code{POST /api/hpc/inbox/import} starts a local job of
kind \code{import} that calls \code{Working.training.hpc_import.import_results}; \code{GET
/api/hpc/inbox} lists the attempts. Seed results: \code{POST /api/hpc/seed/import}. Review:
\code{GET} of the open queues through \code{webui/client/src/api/review.ts}.
\end{underthehood}

\section{Sending a long run to the cluster}\label{sec:jobs-sending}
\subsection{What runs here and what goes to the cluster}
Before a run starts, the app estimates how long it will take on this machine, from each
\term{block}{block}'s own cost model, and compares that with a \emph{local limit} set per workspace
in Settings \textrightarrow{} Compute \& HPC (\apppage{settings.compute-hpc}). Without a saved
value the limit is 15 minutes in Analyse (\code{CLUSTER_ROUTING_CEILING_S}), 20 minutes for a
Discovery template run (\code{DEFAULT_CEILING_MIN}), and 10 minutes for a seed search
(\code{SEED_CEILING_S}). A run under the limit starts here. A run over it does not start:
\begin{itemize}
\item In Analyse, pressing \ui{Run} on a chain over the limit writes a SLURM job for the chain
  instead, and the message says how long the run would take, which stage is slow, the longest span
  that would run here, and the \code{sbatch} command.
\item On Discovery's Seed page, a seed search over the limit is written as a cluster job
  (\apppage{discovery.seed}).
\item Models \textrightarrow{} Launch writes a CNN training job, and the shape-clustering block page in
  Analyse writes a Ward clustering over every training window
  (\apppage{models.launch}).
\end{itemize}
A run whose cost model is still being timed on this machine is held until the estimate exists, so it
is not started blind.

\begin{plainterms}[What SLURM is]
A \term{cluster}{cluster-hpc} is many computers shared by many people. SLURM is the program that
queues their work: you hand it a short script saying what to run and for how long, with the
\code{sbatch} command, and it runs the script when a computer is free.
\end{plainterms}

\subsection{What is sent, and how results come back}
Every job is a folder or a pair of files under \code{HPC/} in the repository: the recipe or
specification, the \code{sbatch} script, and, in the script's comments, the code, input files and
outputs the job needs. Recording samples travel as their channel files. The cluster runs the same
core code as the laptop, so a cluster result and a local one are the same calculation.

\begin{watchout}[The cluster's 20-minute limit]
This account's cluster limit stops any job after 20 minutes. A seed search saves its progress after
every channel and resubmits itself until it is done, so it survives the limit. A chain written from
Analyse cannot resume: if its estimate is over 20 minutes, the message says it will most likely be
killed and gives the span that would fit. Shorten the span, or split it.
\end{watchout}

Results come back three ways, and each row on the Jobs page says which applies:
\begin{description}
\item[Manifest inbox] A job folder (CNN training, Ward clustering) writes its results into its
  \code{out/} folder. Copy that back and import it. The import checks that the recipe the results
  were made from is the recipe the job was written with, that the windows and the pool are the ones
  in this database, and that no test score would change under a frozen cut. If any check fails,
  nothing is recorded.
\item[Seed import] A seed search writes one small result file. Copy it back and import it from its
  row; it is checked against the search it was written for.
\item[Not through the site] Discovery's chain export and Models' paired training write their runs
  into the cluster's copy of the database. There is no return path for them yet; the row says so.
\end{description}
Importing the same job twice records it once. \Cref{sec:architecture-long-work} gives the code.

\subsection{What ``waiting in review'' means}
Some work is finished by the machine but not by you. When a run's detections, a training set's
windows or a family's members are sent to Review from their own page, they become a review queue.
That queue is listed under \emph{Waiting in review} until every item has a
\term{verdict}{verdict}. Nothing runs while it waits; the board only counts what is left and how fast
you have been going. \Cref{ch:review} covers judging.
